% Problem record op_c37650bfb81dbfc6. This TeX file is derived from
% database/problems_json/op_c37650bfb81dbfc6.json by scripts/sync-tex.mjs; edit the JSON record
% and rerun the script rather than editing this file.
\section{Minimum LU--LC counterexample for graph states}
\label{sec:op_c37650bfb81dbfc6}

\paragraph{Problem.}
What is the least number of qubits for which two graph states can be locally
unitary equivalent without being locally Clifford equivalent?  For a simple
graph $G=(V,E)$ with $V=\{1,\ldots,n\}$, define its graph state by
\begin{equation}
  \lvert G\rangle
  :=\left(\prod_{\{u,v\}\in E}\mathrm{CZ}_{uv}\right)
    \lvert+\rangle^{\otimes n},
  \qquad
  \lvert+\rangle:=\frac{\lvert0\rangle+\lvert1\rangle}{\sqrt2}.
  \label{eq:p48-graph-state}
\end{equation}
For graphs $G$ and $H$ on $n$ vertices, use the state convention in
Eq.~\eqref{eq:p48-graph-state} and write
\begin{equation}
  \begin{aligned}
    G\sim_{\mathrm{LU}}H
    &\iff
      \lvert H\rangle=e^{i\phi}
      \left(\bigotimes_{j=1}^{n}U_j\right)\lvert G\rangle
      &&\text{for some }U_j\in U(2),\ \phi\in\mathbb R,\\
    G\sim_{\mathrm{LC}}H
    &\iff
      \lvert H\rangle=e^{i\theta}
      \left(\bigotimes_{j=1}^{n}C_j\right)\lvert G\rangle
      &&\text{for some }C_j\in\mathcal C_1,\ \theta\in\mathbb R,
  \end{aligned}
  \label{eq:p48-lu-lc-equivalence}
\end{equation}
where $\mathcal C_1$ is the single-qubit Clifford group.  Since
$\mathcal C_1\subset U(2)$, LC equivalence implies LU equivalence.  Define the
minimum counterexample size by
\begin{equation}
  n_{\min}
  :=\min\left\{n:\text{there exist $n$-vertex graphs $G,H$ with
  $G\sim_{\mathrm{LU}}H$ and $G\not\sim_{\mathrm{LC}}H$}\right\}.
  \label{eq:p48-minimum-counterexample}
\end{equation}
Determine the integer in Eq.~\eqref{eq:p48-minimum-counterexample}.

\subsection*{Status}

Solved

\subsection*{Source}

Kr\"uger and Werner record the original conjecture that LU equivalence of
graph states always implies LC equivalence
\sourcecite{ref:p48-krueger-werner}{KW05}.  After Ji, Chen, Wei, and Ying
constructed a 27-qubit counterexample, Claudet isolated and resolved the
minimum-size question in Eq.~\eqref{eq:p48-minimum-counterexample}
\sourcecite{ref:p48-ji-et-al}{JCWY10},
\sourcecite{ref:p48-claudet}{Cla26}.

\subsection*{Progress}

\begin{itemize}
  \item Kr\"uger and Werner posed the universal implication
  $G\sim_{\mathrm{LU}}H\Rightarrow G\sim_{\mathrm{LC}}H$.  This formulation
  did not include a minimum counterexample size
  \sourcecite{ref:p48-krueger-werner}{KW05}.

  \item Ji, Chen, Wei, and Ying disproved the universal implication by
  constructing LU-equivalent but non-LC-equivalent graph states on 27 qubits.
  Their construction proves $n_{\min}\leq27$ but does not by itself exclude
  smaller counterexamples \sourcecite{ref:p48-ji-et-al}{JCWY10}.

  \item Claudet proved that LU and LC equivalence coincide for every pair of
  graph states on at most 26 qubits.  Combining this lower bound with the
  27-qubit construction gives
  \begin{equation}
    n_{\min}=27.
    \label{eq:p48-minimum-resolution}
  \end{equation}
  Thus Eq.~\eqref{eq:p48-minimum-resolution} resolves the minimum-size problem
  \sourcecite{ref:p48-claudet}{Cla26}.
\end{itemize}

\subsection*{References}

\begin{enumerate}[label={},leftmargin=4.8em,labelsep=0.5em]
  \footnotesize
  \item[\textup{[KW05]}]\label{ref:p48-krueger-werner}
  O. Kr\"uger and R. F. Werner (eds.),
  ``Some Open Problems in Quantum Information Theory,''
  arXiv:quant-ph/0504166 (2005), Problem~28, pp.~70--71.
  \href{https://doi.org/10.48550/arXiv.quant-ph/0504166}{doi:10.48550/arXiv.quant-ph/0504166};
  \href{https://arxiv.org/abs/quant-ph/0504166}{arXiv:quant-ph/0504166}.

  \item[\textup{[JCWY10]}]\label{ref:p48-ji-et-al}
  Z. Ji, J. Chen, Z. Wei, and M. Ying,
  ``The LU-LC Conjecture Is False,''
  \emph{Quantum Information and Computation} \textbf{10}, 97--108 (2010).
  \href{https://doi.org/10.26421/QIC10.1-2-8}{doi:10.26421/QIC10.1-2-8};
  \href{https://arxiv.org/abs/0709.1266}{arXiv:0709.1266}.

  \item[\textup{[Cla26]}]\label{ref:p48-claudet}
  N. Claudet, ``The 27-qubit Counterexample to the LU-LC Conjecture Is
  Minimal,'' arXiv:2603.25219v1 (2026).\newline
  \href{https://doi.org/10.48550/arXiv.2603.25219}{doi:10.48550/arXiv.2603.25219};
  \href{https://arxiv.org/abs/2603.25219}{arXiv:2603.25219}.
\end{enumerate}

\subsection*{Comment}

Equation~\eqref{eq:p48-minimum-resolution} concerns the minimum size at which
LU and LC equivalence can differ; it does not say that 27 qubits are required
for LU--LC equivalence.  The original universal conjecture was already
resolved negatively by the 27-qubit construction.  As of September 2026, the
matching lower bound through 26 qubits is contained in a recent arXiv v1
preprint, so the archived solved status records its theorem rather than its
peer-review history.

\subsection*{Field}

Quantum Resource Theory

\subsection*{Topic}

Local unitary equivalence

\subsection*{ID}

\texttt{op\_c37650bfb81dbfc6}
